% Author: Lisa Lusinga Namugenyi (editor). Draft written by Charis Opol on 23 Nov 2026.
% IEEE Access rules (template access.tex): one paragraph, 150-250 words, self-contained, no
% abbreviations, footnotes, references, equations or tables. Every number also appears in
% Sections 08 and 09.
\begin{abstract}
Automatic speech recognition degrades sharply in acoustic noise, which limits its use for live
captioning. This paper evaluates whether visual information from the speaker's lips recovers that
loss, first on the GRID corpus with added noise and then on open-vocabulary speech from films and
TED talks. The proposed model combines an audio encoder and a residual-network lip encoder through
bidirectional cross-modal attention, a per-frame fusion gate, a learned embedding for a missing
stream, and modality dropout. On four
unseen speakers it reduces the word error rate of an audio-only baseline by 44 to 51 percent at
every noise level, from 1.96 to 1.02 percent on clean audio and from 9.45 to 5.34 percent at
a signal-to-noise ratio of \textminus5~dB. Ablations attribute most of this gain to modality
dropout and to the three-dimensional convolutional front end. For real speech, a
second model fuses pretrained self-supervised audio and lip-reading encoders through a
reliability-gated fusion that is initialised as its audio-only twin. On 67 held-out TED talks at
\textminus10~dB it reduces the word error rate from 42.3 to 28.5 percent, at a cost of 1.3
percentage points on clean speech. When the audio does not belong to the visible speaker, the model
transcribes the wrong voice, and retraining on such audio removes this error only at a cost of 15.9
points on clean speech. A shared inference engine serves a command-line tool, a browser
demonstration, and a live capture mode.
\end{abstract}

\begin{keywords}
Audio-visual speech recognition, connectionist temporal classification, cross-modal attention,
lip reading, live captioning, multimodal fusion, noise robustness.
\end{keywords}
